%==============================================================================
%  WEEKLY AI BRIEF -- issue #07, week ending 6 October 2026
%  Engine: tectonic
%==============================================================================
\documentclass[11pt,a4paper]{article}

%------------------------------ packages -------------------------------------
\usepackage[T1]{fontenc}
\usepackage[utf8]{inputenc}
\usepackage{lmodern}
\usepackage[margin=2.2cm,top=2.0cm,bottom=2.3cm]{geometry}
\usepackage{microtype}
\usepackage{graphicx}
\usepackage{xcolor}
\usepackage{titlesec}
\usepackage{enumitem}
\usepackage{booktabs}
\usepackage{tabularx}
\usepackage{fancyhdr}
\usepackage{lastpage}
\usepackage{tcolorbox}
\usepackage{xurl}
\usepackage[hidelinks]{hyperref}

\tcbuselibrary{skins,breakable}

\graphicspath{{../../../template/}{./}{template/}}

%------------------------------ palette --------------------------------------
\definecolor{brandink}{HTML}{103868}    % KAUST navy -- headings, rules
\definecolor{brandaccent}{HTML}{0B7480} % KAUST teal (darkened to 5.6:1 on white)
\definecolor{brandsoft}{HTML}{EAF3F5}   % box fills
\definecolor{brandgrey}{HTML}{5A6672}   % meta text

\hypersetup{colorlinks=true,linkcolor=brandaccent,urlcolor=brandaccent,
            citecolor=brandaccent,pdfborder={0 0 0}}

%------------------------------ metadata -------------------------------------
\newcommand{\briefnumber}{07}
\newcommand{\briefweek}{30 September -- 6 October 2026}
\newcommand{\briefdate}{6 October 2026}
\newcommand{\briefauthor}{Ali Habibullah}
\newcommand{\briefaffil}{KAUST}
\newcommand{\briefscope}{Papers, model releases, benchmarks and policy from the past seven days}

%------------------------------ section styling ------------------------------
\titleformat{\section}
  {\normalfont\sffamily\large\bfseries\color{brandink}}{\thesection.}{0.6em}{}
  [\vspace{-0.55em}{\color{brandaccent}\rule{\textwidth}{0.9pt}}]
\titlespacing*{\section}{0pt}{1.5ex plus .3ex}{0.9ex}

\titleformat{\subsection}[runin]
  {\normalfont\sffamily\bfseries\color{brandink}}{}{0pt}{}[\;---]
\titlespacing*{\subsection}{0pt}{1.25ex plus .2ex}{0.5em}

\setlist[itemize]{leftmargin=1.15em,itemsep=0.25em,topsep=0.35em,parsep=0pt}
\setlength{\parindent}{0pt}
\setlength{\parskip}{0.45em}

%------------------------------ header / footer ------------------------------
\pagestyle{fancy}
\fancyhf{}
\renewcommand{\headrulewidth}{0pt}
\renewcommand{\footrulewidth}{0.4pt}
\fancyfoot[L]{\footnotesize\color{brandgrey}Weekly AI Brief \#\briefnumber\ \textbullet\ \briefweek}
\fancyfoot[R]{\footnotesize\color{brandgrey}Page \thepage\ of \pageref{LastPage}}

%------------------------------ content macros -------------------------------
\newcommand{\tldr}[1]{%
  \begin{tcolorbox}[breakable,enhanced,colback=brandsoft,colframe=brandaccent,
      boxrule=0.8pt,left=10pt,right=10pt,top=8pt,bottom=8pt,arc=2pt,
      title={\sffamily\bfseries The week in five lines},fonttitle=\color{white}]
    \begin{itemize}[leftmargin=1.1em,itemsep=0.3em,topsep=0pt]#1\end{itemize}
  \end{tcolorbox}}

\newcommand{\entry}[4]{%
  \par\vspace{0.7em}%
  {\sffamily\bfseries\color{brandink}#1}\par
  {\footnotesize\color{brandgrey}#2 \textbullet\ \href{#3}{link} \textbullet\ \srcref{#4}}\par
  \vspace{0.15em}}

\newcommand{\why}[1]{\par{\small\textbf{Why it matters.} #1}\par}

\newcommand{\srcref}[1]{\textcolor{brandaccent}{[S#1]}}

\newcommand{\stat}[2]{{\sffamily\bfseries\color{brandink}#1}~{\footnotesize\color{brandgrey}#2}}

%==============================================================================
%==============================================================================
\begin{document}

%------------------------------ title block ----------------------------------
\noindent
\begin{minipage}[c]{0.48\textwidth}
  \raggedright\includegraphics[height=1.4cm]{KAUST_Academy_Logo.png}
\end{minipage}\hfill
\begin{minipage}[c]{0.48\textwidth}
  \raggedleft\includegraphics[height=1.5cm]{KAUST_Logo.png}
\end{minipage}
\vspace{0.7em}

{\sffamily
 {\color{brandgrey}\footnotesize ISSUE \#\briefnumber\ \textbullet\ \briefdate}\par
 \vspace{0.15em}
 {\LARGE\bfseries\color{brandink}Weekly AI Brief}\par
 \vspace{0.25em}
 {\large\color{brandaccent}\briefweek}\par
 \vspace{0.35em}
 {\footnotesize\color{brandgrey}\briefscope\\
  Prepared by \briefauthor, \briefaffil}\par}
\vspace{0.3em}
{\color{brandink}\rule{\textwidth}{1.4pt}}
\vspace{0.6em}

%------------------------------ 1. executive summary -------------------------
\tldr{
  \item \textbf{Google announced Gemini 4 Argon, but only cyber defenders can use it.} It is priced at \$2/\$10 per million tokens at launch. Artificial Analysis ranks it 8th of 224, not first. \srcref{1}\,\srcref{2}
  \item \textbf{OpenAI's GPT-6.1 Sol costs one fifth of GPT-6 Astra.} OpenAI says it nearly matches Astra on agentic coding. Artificial Analysis puts it one point below Astra. \srcref{3}\,\srcref{5}
  \item \textbf{Aleph Alpha released Kolibri-1 under Apache 2.0.} It is a 78B-total, 3.46B-active German/English mixture-of-experts model with a 1M-token context. \srcref{6}
  \item \textbf{Anthropic says Z.ai's open-weight GLM-5.3 writes working exploits.} Simple tricks got past its safeguards 64--100\% of the time. \srcref{17}
  \item \textbf{California signed 13 more AI bills.} One bars employers from firing or disciplining workers on an AI decision alone. \srcref{18}
}

%------------------------------ 2. models & releases -------------------------
\section{Model and product releases}

\entry{Gemini 4 Argon: Google's new frontier model, limited to cyber defenders for now}{Google DeepMind \textbullet\ Fairwind Program only \textbullet\ 30 September}{https://blog.google/innovation-and-ai/models-and-research/gemini-models/gemini-4-argon/}{1}
For now Argon is available only to vetted cyber defenders in Google's Fairwind Program,
who get it ``without cyber guardrails''. Google says paid API customers and Google AI
Ultra subscribers come next. It gives no date and no API model ID. The output limit rises
to 1M tokens, up from 64K. Launch pricing is \$2 input and \$10 output per million tokens,
with cached input 95\% off. The standard price after the launch period is \$4/\$20.
Google's self-reported numbers include 77.9\% on DeepSWE~v1.1, which Google calls a new
state of the art, 51.3\% on AutomationBench (\#1), 91.7\% on LVBench and a tie for first on
CWE-bench~v1 at 68\%. Google also says Argon leads the Vals Index but quotes no competitor
scores. Artificial Analysis, testing independently, scores Argon (High) 53 on Intelligence
Index v4.3.2, \#8 of 224 models. \srcref{2}
\why{Most teams cannot use Argon yet, so they should plan around next week's models, not its launch-blog numbers.}

\entry{GPT-6.1 Sol: OpenAI says near-Astra results at a fifth of the price}{OpenAI \textbullet\ API, ChatGPT Work, Codex \textbullet\ 29 September (DevDay)}{https://openai.com/index/introducing-gpt-6-1-sol/}{3}
This launched at DevDay on 29 September, the day issue \#06 went out, so it was held over to
this issue. \texttt{gpt-6.1-sol} costs \$2 input, \$0.10 cached input and \$10 output per
million tokens. That is one fifth of GPT-6 Astra's \$10/\$50, and the cached-input price is
half of GPT-6 Sol's. OpenAI's self-reported claims: on DeepSWE~v1.1 it matches Astra and is
6.4 points above GPT-6 Sol's best score. On AutomationBench it is 2.2 points above Opus~5.5 at
medium effort. On the OSWorld~2.0 offline set it is within 2.1 points of Astra. On
Terminal-Bench Science it costs \$5.47 per task, against \$23.80 for Astra. Astra still has
the top score there, at 68.1\%. Artificial Analysis scores Sol (Max) at 52, \#11 of 224. \srcref{5}
DevDay also brought ``dots'', which OpenAI calls always-on agents; computer use in the Agents API;
an Ultrafast tier that OpenAI says generates tokens up to 8$\times$ faster in Codex; and a
Decisions API built on Luna, in limited preview. \srcref{4}
\why{When an agent pipeline spends most of its tokens re-reading the same context, the \$0.10 cached-input price matters more than the headline price.}

\entry{Kolibri-1: an Apache-2.0 German/English MoE from Aleph Alpha}{Aleph Alpha \textbullet\ open weights on Hugging Face \textbullet\ posted 2 October}{https://huggingface.co/Aleph-Alpha/Kolibri-1}{6}
Kolibri has 78B total and 3.46B active parameters (384 experts per layer, 6 routed plus 1 shared).
It was trained on 20T tokens: about 62.5\% English, 23.9\% German and 13.6\% code. Its native
context is 262,144 tokens, and Aleph Alpha says it has validated quality up to 1,048,576.
The FP8 weights need about 78\,GB, which fits on one H200 or B200. Aleph Alpha's own table
(same harness for every model) gives Kolibri an English overall average of 75.5, against
74.7 for Qwen3.5-35B-A3B; the German averages are 70.8 and 69.8. On SWE-bench Verified it
scores 66.4, behind Qwen3.6-35B-A3B's 73.8. Pre-training took 392k B200 GPU-hours, and the
card lists total training compute as 6.4e23 FLOPs.
\why{It is a European open model with a published compute budget and a German-specific tokenizer, and it fits on one GPU.}

\entry{Clef: Cloudflare open-sources models that output decisions, not text}{Cloudflare \textbullet\ Apache 2.0, Workers AI \textbullet\ 30 September--1 October}{https://huggingface.co/Cloudflare/clef}{7}
Clef (fine-tuned from Qwen3.8-27B) and Clef-Flash (from Qwen3.5-9B) take a state and a schema of
typed questions. In one forward pass they return a probability for every allowed answer, with
no free-text output to parse. Cloudflare reports 97.43 macro-F1 on CLINC150+OOS and
94.20 on BANKING77 for Clef. It gives a median latency of 38.8\,ms for Clef-Flash,
against 524.1\,ms for its own earlier Jev model. All numbers are self-reported. \srcref{8}
\why{Routing and classification steps in an agent can move to a model that cannot produce malformed JSON.}

%------------------------------ 3. research ----------------------------------
\section{Research worth reading}

\entry{Sharpening Tax: RL post-training trades coverage for consistency in agents too}{Oh, Zeng, \ldots, Mirhoseini, Li \textbullet\ arXiv:2610.01509 \textbullet\ 1 October}{https://arxiv.org/abs/2610.01509}{9}
The authors give base (pre-trained only) LLMs a light inference harness. These base agents
score far lower on pass@1 than their post-trained versions, but with enough samples they
often solve more distinct tasks (higher pass@K). The cause is that post-training pushes each
task towards ``always solved or never solved''. The paper names this cost the ``Sharpening Tax''
and measures it across 14 base/post-trained pairs from four model families on three agentic
benchmarks (42 cases). The tax shows up in most of them and can be estimated from a few rollouts.
A proposed per-prompt temperature sampler (PTGS) pays less of the tax during RL.
The paper improves both single-shot accuracy and pass@K.
\why{If your agent relies on best-of-$K$ sampling with a verifier, test the base checkpoint as well as the post-trained one.}

\entry{RealCompanion: ten real human--AI relationships, 27,218 messages}{Behnam, Kim, Yang \textbullet\ arXiv:2610.01780 \textbullet\ 1 October}{https://arxiv.org/abs/2610.01780}{10}
The dataset covers ten real relationships between a person and an AI companion, each lasting
up to 120 days. Only 3.4\% of user messages depend on something said earlier. When one does,
the earlier message is a median 2,157 messages back. A recency window finds it 95.9\% of the
time overall but only 2.2\% of the time when it is far back. Relabelling the same history as
``memories'' makes models bring up the past 10--14 points more often, even when it is
irrelevant. The data is released.
\why{Memory evaluations built from invented personas probably overstate how often recall matters, and averaging over all messages hides the long-range cases where it fails.}

\entry{Science or Slop?: detecting AI-written papers by their reasoning, not their tokens}{Oh, Lee, Ahn, Kim, Kang \textbullet\ arXiv:2610.00531 \textbullet\ 30 September}{https://arxiv.org/abs/2610.00531}{11}
SciSlopBench pairs 390 AI-generated papers with human papers matched on research problem and
contribution type. It scores six measures of structure, argument and artifacts. Together they
identify the AI paper in each pair 85.9\% of the time, against 68.7\% for the Binoculars
token-level detector. Higher scores track lower ICLR ratings in every year from 2017 to 2025.
Prompting a model to reduce the measures directly led to reward hacking. The authors'
evidence-grounded revision harness instead closes 63\% of the remaining gap.
\why{Reviewers and program committees get a detector aimed at broken reasoning rather than word choice, which is harder to evade by paraphrasing.}

\entry{Kandinsky 6.0 Video: open text-to-video with synchronized audio, MIT licence}{Team Kandinsky \textbullet\ arXiv:2610.05608 \textbullet\ 4 October}{https://arxiv.org/abs/2610.05608}{12}
The family has two sizes: Lite (3B) and Pro (29B). Both generate 5-second clips with
synchronized 44\,kHz audio, including lip-sync, from text or an image, and a built-in
super-resolution model upscales to 1920$\times$1080. Video and audio streams are linked by
two-way cross-attention. The checkpoints have been appearing on Hugging Face since
9 September; the new item this week is the technical report. The quality claims come from
the team's own side-by-side human evaluation. It reports Pro clearly beats Kandinsky~5.0 and
is ``competitive'' with leading models, with no independent ranking cited.
\why{It is one of the few open-weight models that generate video and speech together, under a licence that allows commercial use.}

%------------------------------ 4. benchmarks & evals ------------------------
\section{Benchmarks, evaluations and reproductions}

Two labs announced frontier models this week, and each claims a lead on its own benchmark
table. Artificial Analysis runs the same tests on every model, and its index shows them
in a tight cluster below Anthropic's September releases:

{\small
\begin{tabularx}{\textwidth}{@{}X r r l@{}}
\toprule
\textbf{Model (AA variant)} & \textbf{AA Index v4.3.2} & \textbf{Rank / 224} & \textbf{Released} \\
\midrule
Claude Opus 5.5            & 58 & \#1  & Sep \srcref{13} \\
Claude Sonnet 5.5          & 56 & \#2  & 28 Sep \srcref{14} \\
Claude Fable 5.1           & 53 & \#5  & Sep \srcref{15} \\
GPT-6 Astra                & 53 & \#7  & Sep \srcref{16} \\
Gemini 4 Argon (High)      & 53 & \#8  & 30 Sep \srcref{2} \\
GPT-6.1 Sol (Max)          & 52 & \#11 & 29 Sep \srcref{5} \\
\bottomrule
\end{tabularx}}
\vspace{0.4em}

\textbf{Is the movement real?} Google's 77.9\% on DeepSWE~v1.1 is self-reported and so far
not reproduced. It sits above the scores reported for Claude Fable~5 and GPT-6 Sol
on the same benchmark (issue \#06). AutomationBench is the bigger warning. Google gives Argon 51.3\% and
ranks it first, while OpenAI's GPT-6.1 Sol page compares against Opus~5.5 and notes that
Claude Fable~5.1's cost figure leaves out fallbacks, which happened on about 40\% of tasks.
Last week's issue showed two labs reporting Claude Opus~5 at 26.9\% and 50.3\% on this same
benchmark. Until the harness is pinned, AutomationBench numbers from different labs should
not be compared directly. Ranks 5--8 above are tied at 53, so their order says little.

%------------------------------ 5. industry & policy -------------------------
\section{Industry, funding and policy}

\begin{itemize}
  \item \textbf{Anthropic: GLM-5.3's safeguards are easy to bypass.} Anthropic's Frontier
  Red Team (29 September, held over from \#06) tested Z.ai's open-weight model. It produced
  end-to-end V8 exploits in 50 of 410 ExploitBench attempts; Claude Mythos Preview managed 56.
  A cover-story prompt got it to engage with overtly harmful attack tasks 64\% of the time,
  and prefilling its thinking raised that to 92\%. An abliterated copy, made for about \$4,400
  of compute, engaged 100\% of the time. Safeguarded Claude models stayed at zero. Anthropic
  is a direct competitor, and these are its own simulated tests; NIST's CAISI, which it cites,
  independently judged GLM-5.3 the most cyber-capable open-weight model. \srcref{17}
  \item \textbf{California enacts 13 more AI bills (30 September).} SB~947 bars employers
  from relying only on AI to discipline or fire workers. SB~951 requires disclosure when a mass
  layoff is caused by AI. AB~1331 and AB~1883 limit workplace surveillance. SB~1000 and
  AB~2713 update the AI Transparency Act's provenance rules. AB~1864 sets customer-screening
  rules for gene-synthesis providers. Newsom also signed an executive order keeping the term
  ``Artificial Intelligence'' in state use. \srcref{18} That answers a 29 September White House
  order telling federal agencies to say ``Super Intelligence'' instead. \srcref{19}
  \item \textbf{Anthropic commits \$100M to train 10,000 ``Frontier Deployed Engineers'' by
  end-2027 (2 October).} The course is a multi-day in-person block followed by a 12-week residency
  leading a Claude project at the engineer's own employer. First cohorts include engineers
  from Accenture, Deloitte, McKinsey, Morgan Stanley and Novo Nordisk. \srcref{20}
\end{itemize}

%------------------------------ 6. what to watch -----------------------------
\section{What to watch next week}

\begin{itemize}
  \item \textbf{OpenAI rollouts it says are ``coming days'' away:} GPT-6.1 Sol Ultrafast and
  broad release of the Decisions API. Check whether either has shipped by 13 October. \srcref{3}\,\srcref{4}
  \item \textbf{Gemini 4 Argon outside Fairwind.} Google promises paid-API and AI Ultra access but
  gives no date. Once access opens, check whether anyone outside Google reproduces the 77.9\% DeepSWE score. \srcref{1}
  \item \textbf{A response to the GLM-5.3 findings.} Watch whether Z.ai answers Anthropic's
  report or changes how it releases models, and whether CAISI publishes more open-weight cyber assessments. \srcref{17}
\end{itemize}

%------------------------------ 7. sources -----------------------------------
\clearpage
\section{Sources}
{\footnotesize
\begin{tabularx}{\textwidth}{@{}l X@{}}
\toprule
\textbf{\#} & \textbf{Reference} \\
\midrule
S1  & Google (K. Kavukcuoglu). \emph{Gemini 4 Argon: our next era of frontier intelligence}. Google blog, 30 Sep 2026. \url{https://blog.google/innovation-and-ai/models-and-research/gemini-models/gemini-4-argon/} \\
S2  & Artificial Analysis. \emph{Gemini 4 Argon (High)}. Live page, accessed 6 Oct 2026. \url{https://artificialanalysis.ai/models/gemini-4-argon} \\
S3  & OpenAI. \emph{Introducing GPT-6.1 Sol}. 29 Sep 2026. \url{https://openai.com/index/introducing-gpt-6-1-sol/} \\
S4  & OpenAI. \emph{DevDay 2026 Recap}. 29 Sep 2026. \url{https://openai.com/index/devday-2026-recap/} \\
S5  & Artificial Analysis. \emph{GPT-6.1 Sol (Max)}. Live page, accessed 6 Oct 2026. \url{https://artificialanalysis.ai/models/gpt-6-1-sol} \\
S6  & Aleph Alpha. \emph{Kolibri-1} model card. Hugging Face, 2 Oct 2026. \url{https://huggingface.co/Aleph-Alpha/Kolibri-1} \\
S7  & Cloudflare. \emph{Clef} model card. Hugging Face, 30 Sep 2026. \url{https://huggingface.co/Cloudflare/clef} \\
S8  & Cloudflare. \emph{Clef decision models}. Cloudflare blog, 1 Oct 2026. \url{https://blog.cloudflare.com/clef-decision-models} \\
S9  & Oh et al. \emph{Sharpening Tax in Post-Training}. arXiv:2610.01509, 1 Oct 2026. \url{https://arxiv.org/abs/2610.01509} \\
S10 & Behnam, Kim, Yang. \emph{RealCompanion}. arXiv:2610.01780, 1 Oct 2026. \url{https://arxiv.org/abs/2610.01780} \\
S11 & Oh et al. \emph{Science or Slop?} arXiv:2610.00531, 30 Sep 2026. \url{https://arxiv.org/abs/2610.00531} \\
S12 & Team Kandinsky. \emph{Kandinsky 6.0 Video}. arXiv:2610.05608, 4 Oct 2026. \url{https://arxiv.org/abs/2610.05608} \\
S13 & Artificial Analysis. \emph{Claude Opus 5.5}. Live page, accessed 6 Oct 2026. \url{https://artificialanalysis.ai/models/claude-opus-5-5} \\
S14 & Artificial Analysis. \emph{Claude Sonnet 5.5}. Live page, accessed 6 Oct 2026. \url{https://artificialanalysis.ai/models/claude-sonnet-5-5} \\
S15 & Artificial Analysis. \emph{Claude Fable 5.1}. Live page, accessed 6 Oct 2026. \url{https://artificialanalysis.ai/models/claude-fable-5-1} \\
S16 & Artificial Analysis. \emph{GPT-6 Astra}. Live page, accessed 6 Oct 2026. \url{https://artificialanalysis.ai/models/gpt-6-astra} \\
S17 & Anthropic Frontier Red Team. \emph{GLM-5.3 and the spread of advanced cyber capabilities}. 29 Sep 2026. \url{https://www.anthropic.com/research/glm-5-3-and-the-spread-of-advanced-cyber-capabilities} \\
S18 & Office of the Governor of California. \emph{California's nation-leading AI framework just got stronger}. 30 Sep 2026. \url{https://www.gov.ca.gov/2026/09/30/californias-nation-leading-ai-framework-just-got-stronger-governor-newsom-signs-more-first-in-the-nation-worker-protections-and-more/} \\
S19 & The White House. \emph{Fact Sheet: President Donald J. Trump Inaugurates the Era of Super Intelligence}. 29 Sep 2026. \url{https://www.whitehouse.gov/fact-sheets/2026/09/fact-sheet-president-donald-j-trump-inaugurates-the-era-of-super-intelligence/} \\
S20 & Anthropic. \emph{Anthropic invests \$100 million to train 10,000 engineers}. 2 Oct 2026. \url{https://www.anthropic.com/news/claude-frontier-academy} \\
\bottomrule
\end{tabularx}}

\vspace{0.8em}
{\footnotesize\color{brandgrey}
Compiled \briefdate. Every figure in this brief was read from the linked primary
source on that date. Corrections: \texttt{\briefauthor}.}

\end{document}
